\section{Setting}
\label{sec:setting}

Throughout this paper we use the setting and sensitivity framework of \citet{boyarsky2026}. This section restates this setting.

The source sample pools $S$ sites. An observation is $O=(Y,A,X,C)$, where
$A\in\{0,1\}$ is treatment, $X$ contains unit-level covariates, and
$C\in\R^p$ is the site-level context vector, constant within a site. The
basis of $C$ is fixed before the analysis and may contain raw site
characteristics and their transformations. Site $s$ carries the profile
$c_s\in\R^p$. The source law $\sP$ mixes the sites, and treatment is
randomized within each site with known probability
$e_a(x,c)=\pr_\sP(A=a\mid X=x,C=c)$ bounded away from zero. A target site is defined by a
pair $(\tP,\ct)$: a covariate law we can sample and a context profile
$\ct\in\R^p$. Treatment and outcome data exist only for the source sites.

The framework rests on a site-level fixed-effects response model.

\begin{assumption}[Site-level fixed effects, \citealp{boyarsky2026}]
\label{ass:model}
For $a\in\{0,1\}$,
\begin{equation}
  \E_\sP[Y(a)\mid X,C]=g_a(X)+C^\top\beta_a,
  \label{eq:response-model}
\end{equation}
where $g_a$ is unrestricted and $\beta_a\in\R^p$ is fixed. Treatment is
randomized within each source site, consistency holds, the same
functions and coefficients govern the target site, and all second
moments used below are finite.
\end{assumption}

Write $\pseudo{a}=\one\{A=a\}Y/e_a(X,C)$ for the inverse-probability
pseudo-outcome, $\muC(x)=\E_\sP[C\mid X=x]$, and
$\Ctil=C-\muC(X)$. Averaging \eqref{eq:response-model} over the target
population expresses the target average treatment effect as
\begin{equation}
  \psi_T=\bT+\ellvec^\top(\beta_1-\beta_0),
  \qquad
  \begin{aligned}
    \bT&=\E_{X\sim \tP}\bigl[\E_\sP[\pseudo{1}-\pseudo{0}\mid X]\bigr],\\
    \ellvec&=\ct-\E_{X\sim \tP}\bigl[\muC(X)\bigr].
  \end{aligned}
  \label{eq:target-coefficients}
\end{equation}
The scalar $\bT$ is the covariate-transported part and is identified
under overlap in $X$. The loading $\ellvec$ multiplies the fixed-effect
contrast, which is the only unknown. Each $\beta_a$ satisfies the normal
equations of a partially linear regression \citep{robinson1988root},
\begin{equation}
  \CovC\,\beta_a=\dmom{a},
  \qquad
  \CovC=\E_\sP\bigl[\cov(C\mid X)\bigr],
  \quad
  \dmom{a}=\E_\sP\bigl[\cov(C,\pseudo{a}\mid X)\bigr].
  \label{eq:normal-equations}
\end{equation}
Because $C$ takes only $S$ values, every residual $\Ctil$ lies in the
span of the centered site profiles,
\begin{equation}
  \Vspan=\operatorname{span}\{c_s-c_1:\ s=2,\ldots,S\},
  \qquad r=\dim\Vspan\leq S-1,
  \label{eq:span}
\end{equation}
so $\range(\CovC)\subseteq\Vspan$ and $\rankop(\CovC)\leq S-1$. When
$p>S-1$ the normal equations leave fixed-effect directions undetermined,
and $\psi_T$ is point identified if and only if
$\ellvec\in\range(\CovC)$.

The key innovation of this method is that the analyst is able to encode substantive knowledge in a restriction set
$\restrset$ for the stacked vector
$\beta=(\beta_0^\top,\beta_1^\top)^\top$ as a polyhedron or a finite union
of polyhedra. \citet{boyarsky2026} provide several such candidate restrictions: magnitude caps, effect ordering, sign restrictions, and
budgets. The moments and the restriction define
\begin{equation}
  \feasset
  =\bigl\{\beta:\ \CovC\beta_a=\dmom{a},\ a\in\{0,1\},
  \ \beta\in\restrset\bigr\},
  \label{eq:feasible-set}
\end{equation}
and the sensitivity interval is $[\vminus,\vplus]$ with
\begin{equation}
  \vminus=\bT+\inf_{\beta\in\feasset}\ellvec^\top(\beta_1-\beta_0),
  \qquad
  \vplus=\bT+\sup_{\beta\in\feasset}\ellvec^\top(\beta_1-\beta_0).
  \label{eq:endpoint-programs}
\end{equation}

In this paper, we additionally define $\width=\vplus-\vminus$ as the width of the identified set at the
target. An infeasible program means the restriction contradicts the
source moments. An unbounded program means the restriction leaves a
direction that matters for the target uncontrolled. When $\width$ is a finite value it gives us a natural tool to characterize the generalizability of a particular set of experiments to a new target site.

Two more objects from \citet{boyarsky2026} are used in this paper. The
first is the reduced coordinates. Given a fixed orthonormal basis $\Bspan\in\R^{p\times
r}$ of $\Vspan$, write $\Om=\Bspan^\top\CovC\Bspan$ and
$\xired{a}=\Bspan^\top\dmom{a}$. The equality system in
\eqref{eq:feasible-set} is equivalent to
$\Om\Bspan^\top\beta_a=\xired{a}$, which has full row rank when the
smallest eigenvalue of $\Om$ is positive. We assume this eigenvalue
condition throughout, so that
$\range(\CovC)=\Vspan$ and $\kernel(\CovC)=\Vspan^\perp$. The second component is the estimator as defined by Algorithm 1. This algorithm first estimates the five inputs
$\Gam=(\bT,\ellvec,\CovC,\dmom{0},\dmom{1})$ with cross-fitted
orthogonal scores, fits $\muC$ through a joint site-probability
classifier so that every fitted residual lies in $\Vspan$ exactly, pools
the scores, and solves each bound program once in the reduced
coordinates. We call it the input-debiased estimator, in contrast to the one-step
estimator also studied in \citet{boyarsky2026}. Every estimate reported in this paper uses it.
